\section{Legal Landscape}
\label{sec:legal}

\subsection{Federal Foreclosure of Both Standards}

\textit{Davis v.\ Bandemer}\footnote{\textit{Davis v.\ Bandemer}, 478 U.S.\ 109 (1986).}
established that federal courts could, in principle, hear partisan gerrymandering
claims under the equal protection clause. \textit{Vieth v.\ Jubelirer}\footnote{
\textit{Vieth v.\ Jubelirer}, 541 U.S.\ 267 (2004).} came within one vote of
overruling \textit{Bandemer} and imposing the position that all partisan gerrymandering
claims were nonjusticiable. Justice Kennedy's concurrence in \textit{Vieth} preserved
the possibility of a manageable standard, but Kennedy could not identify one.
\textit{Rucho v.\ Common Cause}\footnote{\textit{Rucho v.\ Common Cause}, 588 U.S.\
684 (2019).} completed the trajectory: federal courts cannot adjudicate partisan
gerrymandering claims, period.

Critically, \textit{Rucho}'s holding does not prefer proportionality or majoritarianism;
it holds that federal courts cannot impose either. Under \textit{Rucho}, a map
that produces zero proportionality (winner-take-all regardless of vote share) is
no more constitutionally infirm, at the federal level, than a map with perfect
proportionality. The normative choice is left to state legislatures, state courts
interpreting state constitutions, and Congress acting under the Elections Clause.

\subsection{Pennsylvania: A Soft Proportionality Floor}

The Pennsylvania Supreme Court's decision in \textit{League of Women Voters v.\
Commonwealth}\footnote{\textit{League of Women Voters v.\ Commonwealth}, 178 A.3d
737 (Pa.\ 2018).} interpreted Pennsylvania's Art.\ I, \S 5 (``Elections shall be
free and equal'') to prohibit maps that ``subordinate'' all legitimate redistricting
criteria to partisan advantage. The court drew a remedial map itself, applying
population equality and compactness criteria---producing a map with Gallagher Index
substantially below the invalidated enacted map.

The PA Supreme Court's analysis does not adopt a strict proportionality standard.
It holds that maps must not deviate from what neutral criteria would produce ``far
beyond'' the neutral baseline---a comparative standard that the \textsc{bisect}
reference distribution operationalizes directly. The PA decision is best understood
as imposing a ``proportionality floor at the neutral-criteria baseline,'' not a
requirement of strict proportionality.

Article I, Section 5 of the Pennsylvania Constitution reads: ``Elections shall be
free and equal; and no power, civil or military, shall at any time interfere to
prevent the free exercise of the right of suffrage.'' The phrase ``free and equal''
has been interpreted by the PA Supreme Court to require both equal weight for
each vote and freedom from manipulation that systematically distorts the translation
from votes to representation. The Gallagher Index gap between neutral and enacted
maps provides the empirical measure of how far the distortion extends.

\subsection{North Carolina: Free Elections Clause}

Article I, Section 10 of the North Carolina Constitution reads: ``All elections
shall be free.'' In \textit{Harper v.\ Hall}\footnote{\textit{Harper v.\ Hall},
868 S.E.2d 499 (N.C.\ 2022).} (later modified by subsequent decisions), the
NC Supreme Court held that this clause prohibits maps drawn ``with the intent,
and effect, of subordinating the interests of non-supporters in order to entrench
a political party in power.'' The court did not adopt a proportionality standard;
it held that the constitutional violation was the \textit{intent} to gerrymander,
evidenced by the partisan data access and metric-optimization behavior of the legislature.

The \textsc{bisect} comparison is relevant to NC's intent analysis: the enacted
NC map's Gallagher Index ($0.14$) is nearly three times the neutral-criteria
baseline ($0.05$), and the gap cannot be explained by geographic constraints
(which affect both maps equally). The disproportionality above the neutral
baseline is circumstantial evidence of intent to manipulate.

\subsection{The Defendant's Proportionality Defense}

The most common defense against proportionality-based partisan gerrymandering
claims is: ``Single-member districts are not required to produce proportional
outcomes; the Constitution does not mandate proportional representation.''
This defense is legally correct at the federal level under \textit{Rucho}
and as a matter of strict constitutional text in most states.

The expert witness response is not to contest this legal claim but to reframe
the standard. The relevant comparison is not enacted LSq vs.\ $0$ (perfect
proportionality) but enacted LSq vs.\ bisect LSq (neutral-criteria proportionality).
The defense concedes that single-member districts cannot produce strict proportionality;
the \textsc{bisect} reference demonstrates how much disproportionality is inherent
in the structure, so that any additional disproportionality above the baseline is
attributable to mapmaker choices rather than to the electoral structure itself.
